% =============================================================
%  Module Description — Discrete Structures (CS103-F26)
%  Shown on https://wosm.academy as the course page for CS103-F26.
%  Compile with: pdflatex cs103-module-description.tex
%  Write "TODO(owner): ..." where information is still missing;
%  the website highlights it.
% =============================================================
\documentclass[11pt,a4paper]{article}

\usepackage[utf8]{inputenc}
\usepackage[T1]{fontenc}
\usepackage{lmodern}
\usepackage[margin=2.2cm]{geometry}
\usepackage{booktabs}
\usepackage{tabularx}
\usepackage{array}
\usepackage{enumitem}
\usepackage{xcolor}
\usepackage{titlesec}
\usepackage{fancyhdr}
\usepackage[hidelinks]{hyperref}

% ---------- Module details (edit these) ----------------------
% \ModuleCode must equal the course's `code` in src/content/courses/.
\newcommand{\ModuleTitle}{Discrete Structures}
\newcommand{\ModuleCode}{CS103-F26}
\newcommand{\Programme}{Foundation programme}
\newcommand{\Level}{Semester 1}
\newcommand{\Semester}{Fall 2026}
\newcommand{\Credits}{TODO(owner)}
\newcommand{\ModuleLeader}{Dr. Ali Khudiyev}
\newcommand{\Email}{TODO(owner)}
\newcommand{\OfficeHours}{TODO(owner)}
\newcommand{\Prerequisites}{None}

% ---------- Styling ------------------------------------------
\definecolor{accent}{RGB}{0,70,130}
\titleformat{\section}{\large\bfseries\color{accent}}{\thesection.}{0.5em}{}
\titlespacing*{\section}{0pt}{1.2em}{0.5em}
\setlist{itemsep=2pt, topsep=4pt}
\setlength{\parindent}{0pt}
\setlength{\parskip}{4pt}

\pagestyle{fancy}
\fancyhf{}
\lhead{\small \ModuleCode\ -- \ModuleTitle}
\rhead{\small Module Description}
\cfoot{\small \thepage}
\renewcommand{\headrulewidth}{0.4pt}

% =============================================================
\begin{document}

\begin{center}
  {\LARGE\bfseries\color{accent} \ModuleTitle}\\[4pt]
  {\large Module Description}
\end{center}
\vspace{6pt}

% ---------- Summary table (not shown on the website: the course page
% shows these details from the macros above) --------------------------
\begin{tabularx}{\textwidth}{>{\bfseries}l X >{\bfseries}l X}
\toprule
Module code   & \ModuleCode    & Credits      & \Credits \\
Programme     & \Programme     & Level        & \Level \\
Semester      & \Semester      & Prerequisites & \Prerequisites \\
Module leader & \ModuleLeader  & Email        & \href{mailto:\Email}{\Email} \\
Office hours  & \multicolumn{3}{l}{\OfficeHours} \\
\bottomrule
\end{tabularx}

% ---------- Sections (shown on the website) ------------------
\section{Module Overview}
The Discrete Structures is the first lecture on required modules of mathematics in
the curriculum. It covers:
\begin{itemize}
  \item Basic concepts of logic and mathematics (in the first 4--5 weeks)
  \item Cardinality and combinatorics, graph theory, number, group theory,
        propositional and first-order logic (afterwards)
\end{itemize}

The course introduces fundamental concepts and key areas of discrete mathematics
that are relevant to students of computer science.

\section{Aims}
\begin{itemize}
  \item TODO(owner): aims of the module.
\end{itemize}

\section{Learning Outcomes}
After completing this course, you will be able to:
\begin{enumerate}[label=LO\arabic*.]
  \item Understand and apply the basic terminology and methods of discrete
        mathematics, including logic, algebraic structures, and algorithmic
        reasoning.
  \item Solve combinatorial problems using appropriate mathematical approaches.
  \item Model and solve problems using graph theory, and quantitatively evaluate
        the efficiency of algorithms.
\end{enumerate}

\section{Syllabus}
\begin{tabularx}{\textwidth}{c X}
\toprule
\textbf{Unit} & \textbf{Topic} \\
\midrule
1  & Logic: Propositions and Connectives \\
2  & Logic: Quantifiers and Predicates \\
3  & Set Theory: Basic Definitions \\
4  & Set Theory: Venn diagrams, KV maps, Standard Identities \\
5  & Natural numbers, Peano axioms, Axiom of Induction \\
6  & Binary Relations \\
7  & Endorelations, reflexive-transitive closure \\
8  & Equivalence relations, partial orders \\
9  & Partial orders: lattices and well-ordering \\
10 & Functions \\
11 & Cardinality \\
12 & Combinatorics I \\
13 & Combinatorics II: binomial/multinomial theorem, Stirling numbers \\
14 & Graph Theory \\
15 & Trees \\
16 & Eulerian and Hamiltonian cycles \\
17 & Planarity \\
18 & Number theory \\
19 & Groups I \\
20 & Groups II: Cyclic groups \\
21 & Groups III: Generated subgroups, Cayley graphs, homomorphisms \\
22 & Propositional Logic I: Syntax and Semantics \\
23 & Propositional Logic II: Validity and Semantics \\
24 & Propositional Logic III: Normal Forms and Equisatisfiable CNF \\
25 & Propositional Logic IV: SAT, DPLL, clause-set representation \\
26 & Propositional Logic V: Resolution \\
27 & First-order Logic \\
\bottomrule
\end{tabularx}

\section{Teaching and Learning Methods}
The module is self-paced. All materials are provided on Moodle.

\begin{tabularx}{0.6\textwidth}{X r}
\toprule
\textbf{Activity} & \textbf{Hours} \\
\midrule
Lectures            & TODO(owner) \\
Labs / Tutorials    & TODO(owner) \\
Independent study   & TODO(owner) \\
\midrule
\textbf{Total}      & \textbf{TODO(owner)} \\
\bottomrule
\end{tabularx}

\section{Assessment}
Your knowledge will be assessed throughout the course with quizzes and the final
exam. The exam evaluates whether you can correctly use the mathematical terms and
concepts introduced in the lectures, such as sets, relations, logic, graphs, and
other mathematical objects. It also tests whether you can select and apply the
right logical, combinatorial, graph-theoretical, or algebraic methods to solve
specific problems.

TODO(owner): weight of each assessment component and the learning outcomes it covers.

\section{Reading List}
\begin{itemize}
  \item K.\ H.\ Rosen: \textit{Discrete Mathematics and Its Applications},
        McGraw-Hill, 1995.
  \item R.\ L.\ Graham, D.\ E.\ Knuth, O.\ Patashnik: \textit{Concrete Mathematics:
        A Foundation for Computer Science}, Addison-Wesley, 1994.
  \item D.\ Gries, F.\ B.\ Schneider: \textit{A Logical Approach to Discrete Math},
        Springer, 1993.
  \item D.\ L.\ Kreher, D.\ R.\ Stinson: \textit{Combinatorial Algorithms:
        Generation, Enumeration, and Search}, CRC Press, 1999.
  \item S.\ Pemmaraju, S.\ Skiena: \textit{Computational Discrete Mathematics:
        Combinatorics and Graph Theory with Mathematica}, Cambridge University
        Press, 2003.
\end{itemize}

\section{Materials and Licence}
Course materials are published under the
\href{https://creativecommons.org/licenses/by-sa/4.0/}{Creative Commons Attribution-ShareAlike 4.0}
licence unless stated otherwise. See \href{https://wosm.academy/licence/}{wosm.academy/licence} for details.

\section{Policies}
\begin{itemize}
  \item \textbf{Attendance:} TODO(owner)
  \item \textbf{Late submissions:} TODO(owner)
  \item \textbf{Academic integrity:} All submitted work must be your own.
        See the \href{https://wosm.academy/code-of-conduct/}{Code of Conduct}.
\end{itemize}

\end{document}
